% GENERATED FILE. Edit canonical lessons, then run npm run generate:books.
% canonical-source-hash: fnv1a64:090ac7844b84b082
% canonical-lessons: ML-C530-urjjasvalamaya, ML-C530-curucurukku, ML-C530-vasi, ML-C530-bhiru, ML-C530-nanikkuka

\chapter{Brisk, Bold and Shy}
\label{ch:ml-brisk-bold-and-shy}
\hlchaptermodality{530}

\begin{chapteropening}
\textbf{By the end of this chapter:} \emph{I can say energetic, briskness, obstinacy, a coward and to feel shy.}

Complete the last lesson of chapter 530: I can say energetic, briskness, obstinacy, a coward and to feel shy.
\end{chapteropening}

\section[\texorpdfstring{\ml{ഊർജ്ജസ്വലമായ} (ūrjjasvalamāya)}{ഊർജ്ജസ്വലമായ (ūrjjasvalamāya)}]{\ml{ഊർജ്ജസ്വലമായ} (ūrjjasvalamāya) — energetic; lively}

\label{lesson:ML-C530-urjjasvalamaya}



\begin{quote}
Before the new one: say the Malayalam for a bad habit, then the Malayalam for tolerance.
\end{quote}

\subsection*{You'll want to know: \ml{ഊർജ്ജസ്വലമായ}}
\textbf{\ml{ഊർജ്ജസ്വലമായ}} — \emph{ūrjjasvalamāya} — \textquotedblleft{}energetic; lively\textquotedblright{}.

\begin{quote}
\ml{ഊർജ്ജസ്വലമായ} \ml{കുട്ടികൾ} \ml{മൈതാനത്ത്} \ml{ഓടിക്കളിച്ചു}.
\end{quote}

\begin{quote}
\emph{ūrjjasvalamāya kuṭṭikaḷ maitānattŭ ōṭikkaḷiccu.}
\end{quote}

\begin{quote}
The lively children ran about playing on the field.
\end{quote}

From \ml{ഊർജ്ജം}, energy. Note the doubled \ml{ജ്ജ} after \ml{ർ}.

\begin{cousinweb}
\textbf{\ml{ഊർജ്ജസ്വലമായ}} (\emph{ūrjjasvalamāya}) is built from Sanskrit \emph{ūrjasvala}, \textquotedblleft{}powerful, strong\textquotedblright{} (Monier-Williams), with the Malayalam adjective ending -\ml{ആയ}.
\end{cousinweb}

\begin{grammarlens}[title={What you've built}]
One more word for describing what you lived and felt.
\end{grammarlens}

\subsection*{Your turn}
\begin{itemize}
\raggedright
  \item \emph{Say it:} \emph{ūrjjasvalamāya}
  \item \emph{Say it:} \emph{ūrjjasvalamāya}, once more
  \item \emph{Say it:} \emph{ūrjjasvalamāya kuṭṭikaḷ maitānattŭ ōṭikkaḷiccu.} — The lively children ran about playing on the field.
  \item \emph{Recall:} say the Malayalam for a bad habit, then the Malayalam for tolerance, then say \emph{ūrjjasvalamāya} again
\end{itemize}

\begin{tcolorbox}[breakable,colback=teal!4,colframe=teal!35!black,title={Before you move on}]
What does \ml{ഊർജ്ജസ്വലമായ} mean? (\textquotedblleft{}Energetic; lively\textquotedblright{}.) Say it once more.
\end{tcolorbox}
\section[\texorpdfstring{\ml{ചുറുചുറുക്ക്} (cuṟucuṟukkŭ)}{ചുറുചുറുക്ക് (cuṟucuṟukkŭ)}]{\ml{ചുറുചുറുക്ക്} (cuṟucuṟukkŭ) — briskness; liveliness}

\label{lesson:ML-C530-curucurukku}



\begin{quote}
Before the new one: say the Malayalam for tolerance, then the Malayalam for energetic.
\end{quote}

\subsection*{You'll want to know: \ml{ചുറുചുറുക്ക്}}
\textbf{\ml{ചുറുചുറുക്ക്}} — \emph{cuṟucuṟukkŭ} — \textquotedblleft{}briskness; liveliness\textquotedblright{}.

\begin{quote}
\ml{എഴുപത്} \ml{വയസ്സായിട്ടും} \ml{മുത്തച്ഛന്} \ml{നല്ല} \ml{ചുറുചുറുക്കാണ്}.
\end{quote}

\begin{quote}
\emph{eḻupatŭ vayassāyiṭṭuṁ muttacchanŭ nalla cuṟucuṟukkāṇŭ.}
\end{quote}

\begin{quote}
Even at seventy, grandfather is very brisk.
\end{quote}

\ml{ചുറുചുറുക്കുള്ള} means brisk and lively.

\begin{cousinweb}
\textbf{\ml{ചുറുചുറുക്ക്}} (\emph{cuṟucuṟukkŭ}) is native Dravidian: DEDR 2713 lists \ml{ചുറുക്ക്}, \textquotedblleft{}sharpness, quickness, zeal\textquotedblright{}, and the doubled form for alacrity (DEDR 2713).
\end{cousinweb}

\begin{grammarlens}[title={What you've built}]
One more word for describing what you lived and felt.
\end{grammarlens}

\subsection*{Your turn}
\begin{itemize}
\raggedright
  \item \emph{Say it:} \emph{cuṟucuṟukkŭ}
  \item \emph{Say it:} \emph{cuṟucuṟukkŭ}, once more
  \item \emph{Say it:} \emph{eḻupatŭ vayassāyiṭṭuṁ muttacchanŭ nalla cuṟucuṟukkāṇŭ.} — Even at seventy, grandfather is very brisk.
  \item \emph{Recall:} say the Malayalam for tolerance, then the Malayalam for energetic, then say \emph{cuṟucuṟukkŭ} again
\end{itemize}

\begin{tcolorbox}[breakable,colback=teal!4,colframe=teal!35!black,title={Before you move on}]
What does \ml{ചുറുചുറുക്ക്} mean? (\textquotedblleft{}Briskness; liveliness\textquotedblright{}.) Say it once more.
\end{tcolorbox}
\section[\texorpdfstring{\ml{വാശി} (vāśi)}{വാശി (vāśi)}]{\ml{വാശി} (vāśi) — obstinacy; determination}

\label{lesson:ML-C530-vasi}



\begin{quote}
Before the new one: say the Malayalam for energetic, then the Malayalam for briskness.
\end{quote}

\subsection*{You'll want to know: \ml{വാശി}}
\textbf{\ml{വാശി}} — \emph{vāśi} — \textquotedblleft{}obstinacy; determination\textquotedblright{}.

\begin{quote}
\ml{ജയിക്കണമെന്ന} \ml{വാശിയോടെ} \ml{അവർ} \ml{കളിച്ചു}.
\end{quote}

\begin{quote}
\emph{jayikkaṇamenna vāśiyōṭe avar kaḷiccu.}
\end{quote}

\begin{quote}
They played with a fierce will to win.
\end{quote}

Good or bad: stubbornness in a child, determination in a player.

\begin{cousinweb}
\textbf{\ml{വാശി}} (\emph{vāśi}) is an everyday word; no source we can open explains it, so its origin is not settled.
\end{cousinweb}

\begin{grammarlens}[title={What you've built}]
One more word for describing what you lived and felt.
\end{grammarlens}

\subsection*{Your turn}
\begin{itemize}
\raggedright
  \item \emph{Say it:} \emph{vāśi}
  \item \emph{Say it:} \emph{vāśi}, once more
  \item \emph{Say it:} \emph{jayikkaṇamenna vāśiyōṭe avar kaḷiccu.} — They played with a fierce will to win.
  \item \emph{Recall:} say the Malayalam for energetic, then the Malayalam for briskness, then say \emph{vāśi} again
\end{itemize}

\begin{tcolorbox}[breakable,colback=teal!4,colframe=teal!35!black,title={Before you move on}]
What does \ml{വാശി} mean? (\textquotedblleft{}Obstinacy; determination\textquotedblright{}.) Say it once more.
\end{tcolorbox}
\section[\texorpdfstring{\ml{ഭീരു} (bhīru)}{ഭീരു (bhīru)}]{\ml{ഭീരു} (bhīru) — a coward}

\label{lesson:ML-C530-bhiru}



\begin{quote}
Before the new one: say the Malayalam for briskness, then the Malayalam for obstinacy.
\end{quote}

\subsection*{You'll want to know: \ml{ഭീരു}}
\textbf{\ml{ഭീരു}} — \emph{bhīru} — \textquotedblleft{}a coward\textquotedblright{}.

\begin{quote}
\ml{ഞാൻ} \ml{ഒരു} \ml{ഭീരുവല്ല}, \ml{ഞാൻ} \ml{പോകും}.
\end{quote}

\begin{quote}
\emph{ñān oru bhīruvalla, ñān pōkuṁ.}
\end{quote}

\begin{quote}
I am no coward; I will go.
\end{quote}

The opposite is \ml{ധൈര്യമുള്ള}, brave. Before an ending it takes -\ml{വ}-: \ml{ഭീരുവല്ല}.

\begin{cousinweb}
\textbf{\ml{ഭീരു}} (\emph{bhīru}) is borrowed from Sanskrit \emph{bhīru}, \textquotedblleft{}timid\textquotedblright{}, from the root \emph{bhī}, \textquotedblleft{}to fear\textquotedblright{}, the root of \ml{ഭയം} too (CDIAL 9516).
\end{cousinweb}

\begin{grammarlens}[title={What you've built}]
One more word for describing what you lived and felt.
\end{grammarlens}

\subsection*{Your turn}
\begin{itemize}
\raggedright
  \item \emph{Say it:} \emph{bhīru}
  \item \emph{Say it:} \emph{bhīru}, once more
  \item \emph{Say it:} \emph{ñān oru bhīruvalla, ñān pōkuṁ.} — I am no coward; I will go.
  \item \emph{Recall:} say the Malayalam for briskness, then the Malayalam for obstinacy, then say \emph{bhīru} again
\end{itemize}

\begin{tcolorbox}[breakable,colback=teal!4,colframe=teal!35!black,title={Before you move on}]
What does \ml{ഭീരു} mean? (\textquotedblleft{}A coward\textquotedblright{}.) Say it once more.
\end{tcolorbox}
\section[\texorpdfstring{\ml{നാണിക്കുക} (nāṇikkuka)}{നാണിക്കുക (nāṇikkuka)}]{\ml{നാണിക്കുക} (nāṇikkuka) — to feel shy; to blush}

\label{lesson:ML-C530-nanikkuka}



\begin{quote}
Before the new one: say the Malayalam for obstinacy, then the Malayalam for a coward.
\end{quote}

\subsection*{You'll want to know: \ml{നാണിക്കുക}}
\textbf{\ml{നാണിക്കുക}} — \emph{nāṇikkuka} — \textquotedblleft{}to feel shy; to blush\textquotedblright{}.

\begin{quote}
\ml{അതിഥികളെ} \ml{കണ്ട്} \ml{കുട്ടി} \ml{നാണിച്ചു} \ml{നിന്നു}.
\end{quote}

\begin{quote}
\emph{atithikaḷe kaṇṭŭ kuṭṭi nāṇiccu ninnu.}
\end{quote}

\begin{quote}
The child stood shyly when she saw the guests.
\end{quote}

Past: \ml{നാണിച്ചു}, \emph{nāṇiccu}.

The noun \ml{നാണം}, shyness, is already familiar.

\begin{cousinweb}
\textbf{\ml{നാണിക്കുക}} (\emph{nāṇikkuka}) is native Dravidian, the verb beside \ml{നാണം}, \textquotedblleft{}shame, bashfulness\textquotedblright{} (DEDR 3639).
\end{cousinweb}

\begin{grammarlens}[title={What you've built}]
One more word for describing what you lived and felt.
\end{grammarlens}

\subsection*{Your turn}
\begin{itemize}
\raggedright
  \item \emph{Say it:} \emph{nāṇikkuka}
  \item \emph{Say it:} \emph{nāṇikkuka}, once more
  \item \emph{Say it:} \emph{atithikaḷe kaṇṭŭ kuṭṭi nāṇiccu ninnu.} — The child stood shyly when she saw the guests.
  \item \emph{Recall:} say the Malayalam for obstinacy, then the Malayalam for a coward, then say \emph{nāṇikkuka} again
\end{itemize}

\begin{tcolorbox}[breakable,colback=teal!4,colframe=teal!35!black,title={Before you move on}]
What does \ml{നാണിക്കുക} mean? (\textquotedblleft{}To feel shy; to blush\textquotedblright{}.) Say it once more.
\end{tcolorbox}
